% Options for packages loaded elsewhere
\PassOptionsToPackage{unicode}{hyperref}
\PassOptionsToPackage{hyphens}{url}
\PassOptionsToPackage{dvipsnames,svgnames,x11names}{xcolor}
%
\documentclass[
  11pt,
]{article}
\usepackage{amsmath,amssymb}
\usepackage{iftex}
\ifPDFTeX
  \usepackage[T1]{fontenc}
  \usepackage[utf8]{inputenc}
  \usepackage{textcomp} % provide euro and other symbols
\else % if luatex or xetex
  \usepackage{unicode-math} % this also loads fontspec
  \defaultfontfeatures{Scale=MatchLowercase}
  \defaultfontfeatures[\rmfamily]{Ligatures=TeX,Scale=1}
\fi
\usepackage{lmodern}
\ifPDFTeX\else
  % xetex/luatex font selection
\fi
% Use upquote if available, for straight quotes in verbatim environments
\IfFileExists{upquote.sty}{\usepackage{upquote}}{}
\IfFileExists{microtype.sty}{% use microtype if available
  \usepackage[]{microtype}
  \UseMicrotypeSet[protrusion]{basicmath} % disable protrusion for tt fonts
}{}
\makeatletter
\@ifundefined{KOMAClassName}{% if non-KOMA class
  \IfFileExists{parskip.sty}{%
    \usepackage{parskip}
  }{% else
    \setlength{\parindent}{0pt}
    \setlength{\parskip}{6pt plus 2pt minus 1pt}}
}{% if KOMA class
  \KOMAoptions{parskip=half}}
\makeatother
\usepackage{xcolor}
\usepackage[margin=1in]{geometry}
\usepackage{longtable,booktabs,array}
\usepackage{calc} % for calculating minipage widths
% Correct order of tables after \paragraph or \subparagraph
\usepackage{etoolbox}
\makeatletter
\patchcmd\longtable{\par}{\if@noskipsec\mbox{}\fi\par}{}{}
\makeatother
% Allow footnotes in longtable head/foot
\IfFileExists{footnotehyper.sty}{\usepackage{footnotehyper}}{\usepackage{footnote}}
\makesavenoteenv{longtable}
\setlength{\emergencystretch}{3em} % prevent overfull lines
\providecommand{\tightlist}{%
  \setlength{\itemsep}{0pt}\setlength{\parskip}{0pt}}
\setcounter{secnumdepth}{-\maxdimen} % remove section numbering
\ifLuaTeX
  \usepackage{selnolig}  % disable illegal ligatures
\fi
\IfFileExists{bookmark.sty}{\usepackage{bookmark}}{\usepackage{hyperref}}
\IfFileExists{xurl.sty}{\usepackage{xurl}}{} % add URL line breaks if available
\urlstyle{same}
\hypersetup{
  pdfauthor={Stefano Alimonti --- stefalimonti@gmail.com},
  colorlinks=true,
  linkcolor={blue},
  filecolor={Maroon},
  citecolor={Blue},
  urlcolor={blue},
  pdfcreator={LaTeX via pandoc}}

\author{Stefano Alimonti --- stefalimonti@gmail.com}
\date{March 2026}

\begin{document}

{
\hypersetup{linkcolor=black}
\setcounter{tocdepth}{2}
\tableofcontents
}
\section{Reader's Guide: A GOE-GUE Transition in Carry Companion
Matrices}\label{readers-guide-a-goe-gue-transition-in-carry-companion-matrices}

\textbf{Stefano Alimonti} --- March 2026

\emph{This guide introduces the carry companion matrix ensemble for readers
familiar with random matrix theory but not with carry arithmetic. It explains
where the matrices come from, why their correlations are structured, and what
drives the universality-class transition.}

\begin{center}\rule{0.5\linewidth}{0.5pt}\end{center}

\newpage

\subsection{1. Where the Matrices Come From}\label{where-the-matrices-come-from}

When two integers \(p\) and \(q\) are multiplied in base 2, the carry chain
\(c_0, c_1, \ldots, c_D\) records the overflow at each column position (see
{[}A{]} for the spectral theory of carries). The Carry Representation Theorem
{[}B{]} shows that

\[g(x) h(x) = f(x) + (x - 2) Q(x)\]

where \(g, h, f\) are the digit polynomials of \(p\), \(q\), and \(N = pq\), and
\(Q(x)\) is the carry quotient polynomial with coefficients \(q_k = -c_{k+1}\).
After the standard monic normalization/indexing convention for companion
matrices, the \textbf{carry companion matrix} \(M\) is a \(D \times D\) matrix
(where \(D = \deg Q\)) with ones on the sub-diagonal and the carry values
\(c_1, \ldots, c_D\) in the last column (since the entries are
\(-q_k = c_{k+1}\)).

For a single multiplication, \(M\) is deterministic. Over the ensemble of all
\(D\)-bit semiprimes \(N = pq\) (with \(p, q\) drawn uniformly from \(d\)-bit
primes), the carry coefficients become random, producing a natural random matrix
ensemble.

\textbf{Key structural feature:} This is a \emph{sparse} non-Hermitian ensemble
--- randomness lives only in the last column. The remaining entries are fixed
(ones on the sub-diagonal, zeros elsewhere). This places the ensemble outside
the standard Wigner, Wishart, and Ginibre frameworks.

\begin{center}\rule{0.5\linewidth}{0.5pt}\end{center}

\newpage

\subsection{2. The Correlation Structure}\label{the-correlation-structure}

What makes this ensemble unusual is that the random entries in the last column
are \textbf{not independent}. Adjacent carry values share input bits:
\(c_{k+1} = \lfloor(\text{conv}_k + c_k)/2\rfloor\) depends on \(c_k\), which in
turn depends on \(c_{k-1}\). This creates a Markov correlation structure with:

\begin{itemize}
\tightlist
\item
  \textbf{Lag-1 autocorrelation}
  \(\text{Corr}(c_2, c_3) = 1/\sqrt{2} \approx 0.707\) (Proposition 1, proved
  exactly).
\item
  \textbf{Bulk autocorrelation} numerically observed to increase with position
  \(j\) toward 1.
\item
  \textbf{Spectral gap} \(\rho = 1/2\) (the Diaconis-Fulman eigenvalue {[}A{]}),
  governing the exponential decorrelation rate.
\end{itemize}

If the last-column entries were i.i.d. (same marginals, no correlation), the
ensemble would exhibit GUE-like statistics. The Markov correlations then
suppress the complex-like behavior of the i.i.d. model: they increase the
real-eigenvalue fraction, reduce the effective perturbation rank, and push the
statistics toward the GOE side of the interpolation.

\begin{center}\rule{0.5\linewidth}{0.5pt}\end{center}

\newpage

\subsection{3. The Effective Rank Mechanism}\label{the-effective-rank-mechanism}

The transition is quantified by the \textbf{normalized effective rank} of the
covariance matrix. For a stationary Markov chain with lag-1 correlation \(\rho\)
and length \(D\), the covariance matrix is Kac-Murdock-Szego Toeplitz with
entries \(\Sigma_{ij} = \sigma^2 \rho^{|i-j|}\).

\textbf{Proposition 2:} The normalized effective rank (participation ratio) is

\[g(\rho) = \frac{1 - \rho^2}{1 + \rho^2}\]

in the \(D \to \infty\) limit. At the physical value \(\rho = 1/2\)
(Diaconis-Fulman spectral gap):

\[g(1/2) = \frac{3}{5} = 0.60\]

The correlated random vector effectively has only 60\% of the degrees of freedom
of an uncorrelated one. This rank reduction is the mechanism driving the GOE
shift: fewer effective independent directions means more ``real-like'' spectral
statistics.

\begin{center}\rule{0.5\linewidth}{0.5pt}\end{center}

\newpage

\subsection{4. The Transition in Four
Statistics}\label{the-transition-in-four-statistics}

The GOE-GUE transition is confirmed independently by four spectral observables:

\begin{longtable}[]{@{}
  >{\raggedright\arraybackslash}p{(\columnwidth - 6\tabcolsep) * \real{0.2500}}
  >{\raggedright\arraybackslash}p{(\columnwidth - 6\tabcolsep) * \real{0.2500}}
  >{\raggedright\arraybackslash}p{(\columnwidth - 6\tabcolsep) * \real{0.2500}}
  >{\raggedright\arraybackslash}p{(\columnwidth - 6\tabcolsep) * \real{0.2500}}@{}}
\toprule\noalign{}
\begin{minipage}[b]{\linewidth}\raggedright
Statistic
\end{minipage} & \begin{minipage}[b]{\linewidth}\raggedright
GUE prediction (\(\beta = 2\))
\end{minipage} & \begin{minipage}[b]{\linewidth}\raggedright
Carry matrices
\end{minipage} & \begin{minipage}[b]{\linewidth}\raggedright
GOE prediction (\(\beta = 1\))
\end{minipage} \\
\midrule\noalign{}
\endhead
\bottomrule\noalign{}
\endlastfoot
Spacing ratio \(\langle\tilde{r}\rangle\) & 0.603 & 0.50-0.58 & 0.536 \\
\(\beta_{\text{eff}}\) & 2.0 & \(\leq 2 - \varepsilon\) & 1.0 \\
Real eigenvalue fraction & low & increasing with \(\rho\) & high \\
Number variance \(\Sigma^2(L)\) & GUE curve & intermediate & GOE curve \\
\end{longtable}

The interpolation is continuous: as the correlation parameter \(\lambda\) varies
from 0 (i.i.d.) to 1 (physical Markov), \(\beta_{\text{eff}}\) decreases
monotonically from \(\approx 1.8\) (near-GUE) toward the GOE side, with a slight
sub-GOE endpoint in the largest-\(D\) physical data. This is analogous to the
\(\beta\)-ensembles of Dumitriu and Edelman, but driven by a \emph{physical}
Markov correlation rather than an artificial interpolation.

\begin{center}\rule{0.5\linewidth}{0.5pt}\end{center}

\newpage

\subsection{5. Important Caveats}\label{important-caveats}

\begin{itemize}
\tightlist
\item
  All spectral statistics are \textbf{empirical} (high confidence, no rigorous
  proof).
\item
  Carry matrices are \emph{sparse companion matrices with integer entries} ---
  they do not belong to any standard RMT ensemble.
\item
  The parallel with Riemann zero statistics (GUE) is suggestive but not formal:
  carry matrices approximate Euler factors of \(\zeta(s)\) {[}B{]}, but the
  connection to zero statistics is indirect.
\item
  The GOE-GUE transition is driven by the Markov structure of \emph{carries in
  multiplication}, which is fundamentally different from the time-reversal
  symmetry that distinguishes GOE from GUE in physics.
\end{itemize}

\begin{center}\rule{0.5\linewidth}{0.5pt}\end{center}

\newpage

\subsection{6. Connection to the Broader
Project}\label{connection-to-the-broader-project}

The carry companion matrix ensemble connects to several other results:

\begin{itemize}
\tightlist
\item
  \textbf{Spectral theory {[}A{]}:} The Diaconis-Fulman eigenvalue
  \(\rho = 1/2\) that controls the correlation is the same spectral gap that
  governs carry mixing.
\item
  \textbf{Zeta approximation {[}B{]}:} The ensemble-averaged spectral
  determinant \(\langle|\det(I - M/l^s)|\rangle\) approximates
  \(|1 - l^{-s}|^{-1}\), linking each matrix to an Euler factor.
\item
  \textbf{Covariance structure {[}F{]}:} The exact second-moment structure (Cov,
  autocorrelation) used in Propositions 1-2 is proved rigorously in the
  companion paper.
\end{itemize}

\begin{center}\rule{0.5\linewidth}{0.5pt}\end{center}

\newpage

\subsection{References}\label{references}

\begin{enumerate}
\def\labelenumi{\arabic{enumi}.}
\tightlist
\item
  H.L. Montgomery, ``The pair correlation of zeros of the zeta function,''
  \emph{Proc. Symp. Pure Math.} 24, 181--193, 1973.
\item
  N.M. Katz, P. Sarnak, \emph{Random Matrices, Frobenius Eigenvalues, and
  Monodromy}, AMS, 1999.
\item
  {[}A{]} S. Alimonti, ``Spectral Theory of Carries in Positional
  Multiplication,'' this series.
  doi:\href{https://doi.org/10.5281/zenodo.18895593}{10.5281/zenodo.18895593}
  ---
  \href{https://github.com/stefanoalimonti/carry-arithmetic-A-spectral-theory}{GitHub}
\item
  {[}B{]} S. Alimonti, ``Carry Polynomials and the Euler Product,'' this series.
  doi:\href{https://doi.org/10.5281/zenodo.18895597}{10.5281/zenodo.18895597}
  ---
  \href{https://github.com/stefanoalimonti/carry-arithmetic-B-zeta-approximation}{GitHub}
\item
  {[}F{]} S. Alimonti, ``Exact Covariance Structure of Binary Carry Chains,''
  this series.
  doi:\href{https://doi.org/10.5281/zenodo.18895607}{10.5281/zenodo.18895607}
  ---
  \href{https://github.com/stefanoalimonti/carry-arithmetic-F-covariance-structure}{GitHub}
\end{enumerate}

\begin{center}\rule{0.5\linewidth}{0.5pt}\end{center}

\emph{CC BY 4.0}

\end{document}
